Medical vision--language models (VLMs) are routinely summarized with three numbers: how robust they are to image degradation, how much they rely on the image, and how well they transfer across datasets. Each number is usually reported as a property of the model. We show that each is instead a property of the \emph{model--measurement system}---the model's behaviour plus the design choices that make it observable---and that four unstated preconditions decide whether the number means anything at all: the model must produce variable outputs on real images (\emph{responsiveness}), the counterfactual contrast must be compared against a null in the same metric space (\emph{null model}), the operating point must be reported alongside recall (\emph{operating point}), and the answer contract and output mode must be declared (\emph{contract coverage}).

We test this on six open-weight VLMs evaluated on chest X-rays---two CXR/medical and four general, from four families; 4-bit is post-hoc NF4, bf16/fp16 as released---across three layers---in-domain NIH ChestX-ray14 (1,506 studies $\times$ a 19-cell synthetic degradation grid), cross-hospital PadChest (1,500 studies, stored resolution plus blank), and report-generation Open-i (947 studies)---for 270,324 model calls in the six-model matrix (275,618 including the prose-only seventh model) under a single set of frozen decoding and preprocessing rules, with per-model processor geometry declared. Violating each precondition reverses the conclusion a reader would otherwise draw. Without a responsiveness check, five of six models appear ``image-blind'' on PadChest; in fact they emit \textbf{1--20 distinct answers for 1,500 images}, so reliance is unmeasurable, and the collapse persists on the 265 disease-positive studies. Without a null model, Qwen3-VL appears strongly image-dependent on Open-i (Jaccard 0.09); in fact its blank reply is a single refusal string, so the contrast measures absence detection. Without operating-point reporting, contrast degradation appears to \emph{improve} recall (+4.69\,\pp{}); in fact the positive-prediction rate rises +4.06\,\pp{}, precision falls, and false positives grow from 1,646 to 1,905. Without an output-mode declaration, 4-bit quantization appears to suppress reliance 2.21$\times$ in MedGemma-1.5; in fact the model switches decoding regime (reasoning-channel rate 0.13\% $\rightarrow$ 99.07\% on real images), so the two numbers are not on the same scale.

Two positive results survive these controls. First, degradation acts on \textbf{two separable axes}---whether a model-side monitoring marker fires and how much the corruption \emph{damages} evidence---and noise sits in the dangerous quadrant: damage comparable to resolution loss (pooled over the four readable models, $-0.128$ vs.\ $-0.169$ recall change) while the marker fires 156$\times$ less often (0.46\% vs.\ 71.98\%, in the only model of the six-model matrix whose marker channel is readable). Second, the label-level ordering of degradation sensitivity is \textbf{model-invariant on NIH in raw ranks but partly confounded by prevalence} (pairwise Spearman 0.66--0.86; partial Spearman after rank-residualising prevalence ranges from $-0.69$ to $+0.85$ across model pairs), a candidate design rule rather than a demonstrated lesion law. Finally, we report a negative result: composite risk signals change finding-level error ranking only marginally (AUROC within $+0.034$ of confidence alone) and no model covers more than 9.3\% of findings at a 5\% risk budget.

Confirmatory experiments close the three single-checkpoint gaps the original design left open, on the same 1,506 studies and 31,626 additional cells (calls are higher where a budget was re-run); the schema result now replicates in a second family (Qwen3-VL), while the other confirmatory claims remain single-family. A supervised DenseNet-121 trained only on disjoint patients reaches macro AUROC \textbf{0.733} where the six VLMs sit at 0.50--0.63, so chance-level prompted discrimination is not a task artefact; a frozen-encoder linear probe reaches 0.651--0.703, locating the loss in the prompted readout rather than the image representation; at a fixed input edge the clean and contrast AUROC ordering of two models is stable (0.601--0.611 against 0.520--0.550), but their noise-responsiveness ordering reverses between 512 and 2,048 pixels, so input geometry is itself part of the measurement; the same CNN, run over the 19 cells with a fixed 0.5 rule, also moves its operating point under contrast (+2.1\,\pp{} recall), so Reversal~3 is not specific to vision-language models. A five-schema sweep on a second medical model reproduces the blank-image schema dependence (pathology recall 0.000 in all five schemas) while showing that \emph{which} failure form each schema takes is model-specific. A paraphrase control reproduces the real-image result on MedGemma-4B ($\kappa$ 0.62--0.64) and on Qwen3-VL ($\kappa$ 0.70--0.72), but blank-image answers respond to wording in a model-specific way: p1 flips MedGemma's blank answer and leaves Qwen3's intact. A second reasoning model (Qwen3-VL-8B-Thinking) reproduces the operating-point mechanism---positive-call rate 0.250 $\rightarrow$ 0.125 with macro AUROC flat at 0.53--0.55---and exposes a failure mode the first design could not see: only \textbf{89.0\%} of its cells produce a usable answer at a 5,120-token budget (89.2\% after an 8,192-token sample), coverage itself falls with damage (92.6\% $\rightarrow$ 80.9\%), and the replies that fail expand to fill every budget we tried (median 9.8k, 13.5k, 19.3k, 28.1k characters at caps of 2,560, 3,584, 5,120, 8,192). The axis pair generalises across the two models (noticeability--damage Spearman 0.49 in both, $n=6$, $p>0.26$) but the quadrants do not: noise is silent for MedGemma-1.5's marker channel and loud for a channel that watches termination. Silent failure is a property of the monitor, not of the corruption.

We do not propose a new model. We propose that reliance, robustness, and transfer be reported as \emph{identifiable} quantities---in the audit sense of Section~\ref{sec:preconditions}, not as necessary-and-sufficient conditions---and we give the four checks, the artifacts, and the failure taxonomy needed to do so. A further 9,400-call probe of a frontier-scale open-weight model, served through an OpenAI-compatible relay and run through the same operators, prompts and parser, reproduces the mechanism at larger scale: it emits \textbf{three distinct answers for 500 cross-hospital studies} while emitting 92 for 400 in-domain studies, its null separation is positive but collapses to 0.014 on that layer, and contrast degradation raises both recall and the positive-call rate on negative studies. A 32,184-call cross-modal panel (1,788 brain MRI, retinal OCT and dermoscopy images; three local open-weight checkpoints) makes the mechanism a four-model result: dermoscopy is the most responsive modality for every model, while the null separation falls below zero on retinal OCT for MedGemma-1.5 after a 12-test Bonferroni correction, and only nominally for Qwen3-VL on MRI. A 151-call budget repair recovers the 146 dermoscopy cells truncated at 1,024 tokens, and the study delivers 326,830 unique cells (351,329 requests including budget re-runs and rejected attempts). Threshold, pairing, label-mapping, cluster-bootstrap, and multiplicity sensitivity analyses leave the four reversals intact; a seventh model run only on the prose layers is reported as a measured exclusion, and two marginal collapse cells are reported as marginal rather than used as evidence.
